% Principal GDZ PPN236018647,native484lower–488upper/printed477lower–481upper;HKUST485–489.
\begin{center}[2.]\end{center}
\indent Sind \(n\) Veränderliche vorhanden, so können nicht mehr als \(n\) Bedingungsungleichheiten wirksam sein, und dieselben als Gleichungen betrachtet geben unmittelbar die Minimalwerthe, wenn sie alle wirksam sind.

\indent Nennen wir den analytischen Ort für den Werth der Funktion von \(n\) Variabeln: einen Punkt, wenn alle \(n\) einen bestimmten Werth haben, eine Linie für \(n-1\), eine Fläche für \(n-2\), einen Raum für \(n-3\) etc., so ist in unserer Aufgabe der Punkt zu suchen, welcher einer bestimmten Combination der \(n\) Variabeln entspricht. Der Punkt ist also \textls[120]{gefunden}, wenn \(n\) Ungleichheiten wirksam sind, er ist auf einer Linie zu suchen, wenn \(n-1\) wirksam sind, auf einer Fläche, wenn \(n-2\) etc.

\indent Die Funktion, welche ein Minimum werden soll, ist die Summe der Quadrate sämmtlicher Coordinaten oder das Quadrat der Entfernung des gesuchten Punktes vom Ursprung, das Quadrat des radius vector im weiteren Sinn genommen. Sind gar keine Ungleichheiten gegeben, so ist der Ursprung selbst der gesuchte Punkt. Ist Eine wirksame Ungleichheit gegeben, so liegt der Punkt in einem Mannigfaltigen von \(n-1\) Dimensionen, und zwar so, dass er dem Ursprung möglichst nahe liegt; für 2 Wirksame in einer Mannigfaltigkeit von \(n-2\) Dimensionen etc. Jede hinzukommende wirksame Ungleichheit entfernt den Punkt vom Ursprung und beschränkt seinen Spielraum zugleich um Eine Dimension. Die ganze Schwierigkeit liegt im Aufsuchen der wirksamen Ungleichungen.

\indent [Es soll]
\[
x_1^2+x_2^2+\mathinner{\cdotp\cdotp\cdotp}+x_n^2=\mathrm{Min[imum]}
\]
\noindent [sein, und] es seien [die Ungleichungen] gegeben:
\clearpage\SourceLeaf{485}\ArithmeticHead{478}{L}{GEOMETRIE. NACHLASS.}
\[
\begin{array}{@{}l@{\qquad}l@{}}
a_1'x_1+a_2'x_2+\mathinner{\cdotp\cdotp\cdotp}+b'=u';&u'\geqq0,\\
a_1''x_1+a_2''x_2+\mathinner{\cdotp\cdotp\cdotp}+b''=u'';&u''\geqq0,\\
\multicolumn{2}{c}{.\quad .\quad .\quad .\quad .\quad .\quad .\quad .\quad .\quad .\quad .\quad .\quad .\quad .\quad .}\\
a_1^{(m)}x_1+a_2^{(m)}x_2+\mathinner{\cdotp\cdotp\cdotp}+b^{(m)}=u^{(m)};&u^{(m)}\geqq0,
\end{array}
\]
\noindent wobei wir voraussetzen, dass alle linearen Gleichungen \(u',u'',.\ .\ .\) auf die Form gebracht sind, dass die Summe der Quadrate ihrer Coefficienten \(=1\) ist. Dann ist das Verfahren folgendes:

\indent Man sucht zuerst ein System von Werthen der \(x_1,x_2,.\ .\ .\), welches \(n\) von den Ungleichungen zu 0 macht und die übrigen \(m-n\) positiv. Dann bekommt man \(x_1,x_2,.\ .\ .\) als Funktionen von \(u',.\ .\ .,u^{(n)}\) und zwar, worin letztere \(=0\) gesetzt sind. Substituirt man diese Funktionen in \(u^{(n+1)},.\ .\ .,u^{(m)}\), so hat man
\[
\begin{array}{@{}l@{\quad}l@{}}
(1)&\left\{\begin{array}{@{}r@{}l@{}}
u^{(n+1)}={}&A_1'u'+A_2'u''+\mathinner{\cdotp\cdotp\cdotp}+A_n'u^{(n)}+C',\\
u^{(n+2)}={}&A_1''u'+A_2''u''+\mathinner{\cdotp\cdotp\cdotp}+A_n''u^{(n)}+C'',\\
\multicolumn{2}{c}{.\; .\; .\; .\; .\; .\; .\; .\; .\; .\; .\; .\; .\; .\; .\; .\; .\; .}\\
u^{(m)}={}&A_1^{(m-n)}u'+A_2^{(m-n)}u''+\mathinner{\cdotp\cdotp\cdotp}+A_n^{(m-n)}u^{(n)}+C^{(m-n)},
\end{array}\right.
\end{array}
\]
\noindent worin also, nachdem die Werthe \(u',.\ .\ .,u^{(n)}\) gleich 0 gesetzt sind, nur die Constanten \(C',.\ .\ .,C^{(m-n)}\) bleiben, welche dann nach der obigen Voraussetzung positive Grössen sein sollen. Diesem Werthen-System entspricht dann ein Punkt, welcher in dem von sämmtlichen Ungleichungen abgegrenzten Spielraume liegt, und zwar in der Begrenzung desselben. Den Ungleichungen genügt er demnach, aber im Allgemeinen noch nicht der Bedingung: dass das Quadrat seines Abstandes vom Ursprung der Coordinaten ein Minimum sei. Es kann vielmehr andere Punkte der Begrenzung geben, welche ebenfalls \(n\) von den Ungleichungen gleich 0 und die übrigen positiv machen, dabei aber dem Ursprung näher liegen. Solchen Punkten entsprechen aber Systeme von Werthen der Grössen \(x_1,.\ .\ .,x_n\), welche nicht mehr sämmtliche \(u',.\ .\ .\ u^{(n)}\) zu 0 machen; es werden vielmehr einige von ihnen unwirksam d. h. positiv, und an ihrer Stelle ebensoviele der \(u^{(n+1)},.\ .\ .,u^{(m)}\) zu 0 werden.

\indent Um nun von unserem durch die \(n\) Gleichungen \([u'=0,.\ .\ .,u^{(n)}=0]\) fixirten Ausgangspunkte, den wir als Durchschnittspunkt von \(n\) Mannigfaltigkeiten von \(n-1\) Dimensionen vorstellen, zu solchen näher liegenden Punkten zu gelangen, müssen wir dem Punkte \([x_1,.\ .\ .,x_n]\) eine gewisse Beweglichkeit einräumen, d. h. den Grössen \(u',.\ .\ .,u^{(n)}\), deren Funktion seine Lage ist, eine
\clearpage\SourceLeaf{486}\ArithmeticHead{479}{R}{MANNIGFALTIGKEITEN VON \(n\) DIMENSIONEN.}
\noindent gewisse Veränderlichkeit von ihrem augenblicklichen Werthe 0 aus nach der positiven Seite hin. Um aber gleich die wirksamste Ortsveränderung mit ihm vorzunehmen, d. h. diejenige, mittelst welcher er sich innerhalb der ihm auferlegten Beschränkungen am schnellsten dem Ursprung der Coordinaten nähert, bilden wir zunächst aus den Formeln \(u',u'',.\ .\ .,u^{(n)}\) einen Ausdruck für den radius vector unseres Punktes oder für die Summe
\[
x_1^2+x_2^2+\mathinner{\cdotp\cdotp\cdotp}+x_n^2
\]
\noindent und erhalten daraus
\[
\begin{array}{@{}l@{\qquad}l@{}}
(2)&\left\{\begin{array}{@{}r@{}l@{}}
x_1={}&k_1+\alpha_1'u'+\alpha_1''u''+\mathinner{\cdotp\cdotp\cdotp}+\alpha_1^{(n)}u^{(n)},\\
x_2={}&k_2+\alpha_2'u'+\alpha_2''u''+\mathinner{\cdotp\cdotp\cdotp}+\alpha_2^{(n)}u^{(n)},\\
\multicolumn{2}{c}{.\quad .\quad .\quad .\quad .\quad .\quad .\quad .\quad .\quad .\quad .\quad .}\\
x_n={}&k_n+\alpha_n'u'+\alpha_n''u''+\mathinner{\cdotp\cdotp\cdotp}+\alpha_n^{(n)}u^{(n)}.
\end{array}\right.
\end{array}
\]
\indent Wir untersuchen jetzt, wie die Summe der Quadrate dieser Ausdrücke:
\[
R^2=x_1^2+x_2^2+\mathinner{\cdotp\cdotp\cdotp}+x_n^2
\]
\noindent am schnellsten zu verkleinern ist, d. h. welche von den Grössen \(u',.\ .\ .,u^{(n)}\) wir veränderlich annehmen müssen, um der Variation
\[
\delta(R^2)
\]
\noindent den grössten negativen Werth zu ertheilen. Stellen wir also in der vollständigen Variation die mit \(\delta u',\delta u'',.\ .\ .\) multiplicirten Glieder resp. zusammen, so haben wir nur nach den Factoren derselben zu sehen:
\[
\delta(R^2)=A\delta u'+B\delta u''+\mathinner{\cdotp\cdotp\cdotp}+L\delta u^{(n)}.
\]
\noindent Hat zum Beispiel der Factor \(L\) den grössten negativen Werth, so wissen wir, dass eine Aenderung von \(u^{(n)}\) den Punkt am schnellsten dem Coordinaten-Ursprung nähert. Wir geben also der Grösse \(u^{(n)}\) einen von 0 verschiedenen Werth und lassen denselben so lange zunehmen, bis in Einer der Gleichungen (1) das Glied \(A_nu^{(n)}\) einen ebenso grossen negativen Werth erhält, als das positive Glied \(C\); oder wenn dies nicht der Fall ist, bis der Factor \(L\) durch diese Aenderung aufhört, negativ zu sein, in welchem letztern Falle der gesuchte Punkt gefunden ist. Im ersteren Fall setzt man diejenige der Gleichungen (1), welche 0 wird, an die Stelle von \(u^{(n)}\), welche letztere nunmehr
\clearpage\SourceLeaf{487}\ArithmeticHead{480}{L}{GEOMETRIE. NACHLASS.}
\noindent als unwirksam wegzustreichen ist. Mit den übrigen \(m-1\) Ungleichungen setzt man nun das Verfahren in derselben Weise fort, und zwar so lange, bis in dem Ausdruck \(\delta(R^2)\) sämmtliche Coefficienten positiv.

\indent Mittelst der übrig bleibenden Gleichungen kann man dann ebenso viele \(x,.\ .\ .\) aus der Minimal-Gleichung eliminiren und dann nach den gewöhnlichen Regeln verfahren. Man sucht ein Werthen-System, welches \(n\) von den Ungleichheiten zu 0, die übrigen \(m-n\) positiv macht — wir nennen diese Werthe \(k_1,k_2,.\ .\ .,k_n\) —, und bestimmt dann durch Variation der Grösse \(R^2\) diejenige Aenderung des Werthen-Systems, welche \(R^2\) am schnellsten vermindert, d. h. die kleinste Aenderung des Werthen-Systems, welche eine bestimmte Abnahme der Grösse \(R^2\) um \(\lambda\) bewirkt. Als Maass der Aenderung des Werthen-Systems nehmen wir dabei
\[
\sqrt{\delta x_1^2+\delta x_2^2+\mathinner{\cdotp\cdotp\cdotp}}.
\]
\noindent Wir haben also in dem Ausdruck
\[
2\{x_1\delta x_1+x_2\delta x_2+\mathinner{\cdotp\cdotp\cdotp}\}
\]
\noindent für \(x_1,x_2,.\ .\ .\) die Werthe \(k_1,k_2,.\ .\ .\) zu setzen und ihn dann \(=-[2]\lambda_1\) zu setzen:
\[
\begin{array}{@{}l@{\qquad}l@{}}(3)&k_1\delta x_1+k_2\delta x_2+\mathinner{\cdotp\cdotp\cdotp}=-\lambda_1.\end{array}
\]
\noindent Ferner soll sein
\[
\delta x_1^2+\delta x_2^2+\mathinner{\cdotp\cdotp\cdotp}=\mathrm{Min[imum]}.
\]
\noindent Um die Aenderung des Werthen-Systems \(\delta x_1,\delta x_2,.\ .\ .\) zu finden, welche diese Quadratsumme am schnellsten vermindert, oder die kleinste Aenderung derselben, welche eine bestimmte Verminderung \(\lambda_2\) bewirkt, haben wir wieder
\[
\delta x_1\delta^2x_1+\delta x_2\delta^2x_2+\mathinner{\cdotp\cdotp\cdotp}=-\lambda_2
\]
\noindent und
\[
\delta^2x_1^2+\delta^2x_2^2+\mathinner{\cdotp\cdotp\cdotp}=\mathrm{Min[imum]}.
\]
\noindent Ferner haben wir durch Variation der Gleichung (3):
\[
k_1\delta^2x_1+k_2\delta^2x_2+\mathinner{\cdotp\cdotp\cdotp}=0.
\]
\noindent Wenn wir dann in der Gleichung (3) für \(\delta x_1,.\ .\ .\) ihre Werthe in \(\delta u_1,.\ .\ .\)
\clearpage\SourceLeaf{488}\ArithmeticHead{481}{R}{MANNIGFALTIGKEITEN VON \(n\) DIMENSIONEN.}
\noindent ausgedrückt setzen:
\[
-\lambda_1=K_1\delta u_1+K_2\delta u_2+\mathinner{\cdotp\cdotp\cdotp}+K_n\delta u_n,
\]
\noindent so wird der grösste positive Coefficient \(K_n\) entscheiden, welche Grösse \(u_n\) durch ihre Aenderung zur Verminderung der ursprünglichen Function am meisten beiträgt. Wir ändern deshalb diese, d. h. wir setzen \(u_n>0\) statt \(u_n=0\), und behalten also noch \(n-1\) Gleichungen \(u=0\) übrig, mithin auch
\[
\delta u_1\geqq0,\ \delta u_2\geqq0,\ .\ .\ .,\ \delta u_{n-1}\geqq0.
\]
\noindent Bezeichnen wir \(\delta^2x=\delta(\delta x)\) mit \(\xi\), so soll:
\[
\begin{gathered}
\xi_1\delta x_1+\xi_2\delta x_2+\mathinner{\cdotp\cdotp\cdotp}=-\lambda_2;\\
\xi_1^2+\xi_2^2+\mathinner{\cdotp\cdotp\cdotp}=\mathrm{Min[imum]}
\end{gathered}
\]
\noindent [sein].
